\scanpage{21}
\mspage{\textcircled{2}}

choisissons $u',v'$ \inserted{\unsure{dans|de} $M$} tels que $uu'=vv'$, \inserted{d'où $[u,v]\in G$.} Alors le noyau de l'hom surjectif $G=\pi_1(X_0,e)\longrightarrow\pi_1(X,e)$ est le \inserted{ss-}groupe invariant de $G$ engendré par les éléments suivants, correspondants aux diagrammes (pas nécessairement commutatifs)
\diagram{\begin{tikzcd}[ampersand replacement=\&]x\arrow[r,"f"]\&y\arrow[r,"g"]\&z\\[-1ex]\&e\arrow[ul,"u"]\arrow[u,"v"]\arrow[ur,"w"']\&\end{tikzcd}}
\margin{\uline{NB} S'il existe $\alpha:z\longrightarrow e$, on voit aussitôt que cette relation est triviale. Donc si $e$ est aussi plus grand objet de $X$, alors $\pi_{\unsure{1|0}}(X_0,e)\xrightarrow{\sim}G\xrightarrow{\sim}\pi_1(X,e)$ iso !}
\[
[w,gv][v,fu][w,gfu]^{-1}.
\]

Soient $M$ un monoïde \inserted{\fbox{pseudocofiltrant}$^{*}$}, \inserted{$Y=B_M$,} $\varphi\in M$, considérons la relation d'équivalence \inserted{$R=R_\varphi$} dans $M$
% slanted
\margin{i.e. pour $u,v\in M$, $\exists\,u',v'\in M$ \unsure{avec} $uu'=vv'$}
\[
u\overset{R}{\sim}v\quad\Longleftrightarrow\quad\exists\,p,q\in\mathbb N\qquad\varphi^p u=\varphi^q v.
\]
Cette relation d'équivalence est stable par translations à droite, donc
\[
F\overset{\mathrm{def}}{=}M/\unsure{\mathcal R|R}
\]
est un $M$-ens à droite, i.e un objet de $\deleted{\unsure{X}}Y^\wedge$. On pose $X=Y_{/F}$, donc
% letterhead
\par\smallskip
{\footnotesize\noindent\begin{minipage}[t]{.46\linewidth}\raggedright
\deleted{MINISTERE DE L'EDUCATION NATIONALE}\\ \deleted{ACADEMIE D'AIX-MARSEILLE}\\ \rule{2.6em}{.3pt}\\ \deleted{Inspection Académique}\\ \deleted{de Vaucluse}\\ \rule{2.6em}{.3pt}\\ Tél. : 90 82 99 76\\ \deleted{Référence : Secrétariat n\textsuperscript{o}}
\end{minipage}\hfill%
\begin{minipage}[t]{.51\linewidth}\raggedright
\deleted{AVIGNON, le}\\ \deleted{L'Inspecteur d'académie,}\\ \deleted{directeur des services départementaux}\\ \deleted{de l'Education Nationale de Vaucluse,}\\ \deleted{à M}
\end{minipage}\par}

% archive sheet number: 20
